%% kapa: exam=1 round=21 year=2010 period=2 form=A subject=accounting subject_ko=회계학 seq=30 number=30
%% source: Q-Net 감정평가사 기출문제 (cst00309) artlSeq=5024181 file=extracted/1008/5024181/2010년 제21회 감정평가사 1차시험 기출문제/2교시.hwp
%% conversion: kordoc 4.24.51 -> tools/gen1.py (auto); review=proofread-gpt-sol; flags=none
\begin{question}{30}
20×1년 초 (주)여행은 여객용으로 사용하기 위하여 중고 화물용 항공기를 ￦700에 구입하고, 그 중 일부는 현금을 지급하였다. 여객용으로의 교체는 20×1년 말 완료되었고, 20×2년 초부터 운항에 사용되었다. 다음 자료를 이용하여 아래 물음에 답하시오.
\begin{framedtext}
ㆍ취득 및 등록세 ￦50 \par
ㆍ여객용으로의 교체 원가 ￦100 \par
ㆍ당해 항공기 관련 평균지출액 ￦400 \par
(20×1년 초 지출한 특정차입금 ￦100 포함) \par
ㆍ총차입금 ￦800(특정차입금 ￦100 포함) \par
ㆍ특정차입금 이자율 연 20\% \par
ㆍ차입원가 계산을 위한 일반차입금의 자본화이자율 연 10\% \par
\end{framedtext}
모든 차입금의 차입일은 자본화 개시시점인 20×1년 초이며, 만기는 1년 이상(이자는 연말지급조건)이다. 20×1년 말 항공기의 장부금액은 얼마인가? (단, 특정차입금은 적격자산인 당해 항공기를 취득하기 위한 목적으로 특정하여 차입한 자금을 의미하며, 차입금의 운용으로 인한 수익은 없다고 가정함) \par
\begin{choices}
  \item ￦700
  \item ￦800
  \item ￦900
  \item ￦940
  \item ￦990
\end{choices}
\end{question}
\answer{3}
